\section{信息系统拓扑、服务质量与序贯事件队列}\label{sec:rel-cyber-chronology}

\subsection{信息元件、路径和功能}
信息元件 $q$ 包含固有可用率 $A_q$、单包交付概率 $d_q$、时延 $l_q$、抖动 $j_q$、供电母线、备用能量、功耗和
共因组。功能 $f$ 由一组替代成功路径 $\mathcal P_f$ 描述。路径 $p$ 的确定性服务量为
\begin{equation}
 L_p=\sum_{q\in p}l_q,\qquad J_p=\sum_{q\in p}j_q,\qquad D_p=\prod_{q\in p}d_q.
\end{equation}
只有满足 $L_p\le L_f^{max}$、$J_p\le J_f^{max}$、$D_p\ge D_f^{min}$ 的路径参与功能可用性。零上限表示关闭对应
确定性门限，不表示所有路径时延为零。

\subsection{共享依赖不能按路径独立相乘}
若两条路径共享交换机、电源或控制中心，路径成功事件相关。平台在元件位空间上精确枚举，同一元件只生成一个二值位：
\begin{equation}
 A_f=\sum_{x^c}\mathbf1\!\left[\bigvee_{p\in\mathcal P_f}
 \bigwedge_{q\in p}x_q^c=1\right]p(x^c).
\end{equation}
因此共享元件失效会同时切断所有经过它的路径。离散事件树最多处理二十个二值信息元件，超过限制明确拒绝；结构可靠性
独立接口采用成功路径容斥，最多二十四条约简路径。两个限制针对不同指数维度，不能互相替代。

\subsection{最小割集}
成功路径族的最小击中集就是功能最小割集。平台逐路径扩展候选击中集，并删除包含已有小集合的超集。割集用于解释功能
脆弱点，不直接假定割集互斥。若需要由割集近似不可用度，必须处理割集交叠；当前结构接口直接计算成功路径并集概率，
避免一阶割集和在高不可用率下越过一。

\subsection{联合功能可用性}
FLISR 同时需要检测、隔离和恢复功能，联合成功事件为
\begin{equation}
 \Phi_{DIR}(x)=\phi_D(x)\land\phi_I(x)\land\phi_R(x).
\end{equation}
平台先从每个功能选一条成功路径，再合并所需元件，删除冗余联合路径并做精确容斥。该结果一般不等于
$A_DA_IA_R$。只有三功能依赖元件集合互不相交且条件独立时，标量乘积才成立。

\subsection{共因环境条件化}
天气、站点损毁或控制中心故障以互斥外层环境 $e$ 表示，概率 $p_e$ 和为一。每个环境指定强制失效的共因组和失电母线。
条件于环境后，剩余元件才按独立固有状态枚举：
\begin{equation}
 \Pr(x^c)=\sum_ep_e\Pr(x^c\mid e).
\end{equation}
这比把共因概率分别乘到每条路径更严格，因为同一环境同时改变多个元件。环境概率不归一、未知共因组或非法概率均在
枚举前拒绝。

\subsection{物理供电与备用电池}
信息元件绑定物理母线 $b(q)$。若环境使该母线失电，元件仅在备用自治时间覆盖信息停电时仍可用：
\begin{equation}
 T_q^{aut}=\frac{E_q^{backup}}{P_q^{draw}},\qquad
 x_q^{eff}=x_q^{intr}\mathbf1[T_q^{aut}\ge T^{outage}].
\end{equation}
$P_q^{draw}=0$ 且有备用能量可解释为本模型内不耗能；负功耗或负能量非法。静态离散环境比较持续时间与自治时间，序贯
事件队列则连续更新剩余能量，能表示信息节点在一次长停电中途失效。

\subsection{保护功能绑定}
事件输入把检测、主跳闸、后备跳闸、隔离和恢复分别绑定到信息功能索引。索引 $-1$ 表示本地独立功能，不需要通信；
非负索引必须在功能数组范围内。检测不可用会同时阻断自动隔离与依赖检测的主后备逻辑，恢复只有在隔离成功后才可能成功：
\begin{equation}
 I=D\land\phi_I,qquad R=I\land\phi_R.
\end{equation}
该因果关系避免“未检测故障但恢复成功”的非法类。人工恢复不是把 $R$ 强行置真，而是在三窗口后果中选用人工时长。

\subsection{静态信息状态事件树}
对每个信息状态，保护事件树再展开继电与断路器按需成功分支。联合类概率为
\begin{equation}
 \pi_c=p_e\,p(x^c\mid e)\,p(c_P\mid x^c,e).
\end{equation}
所有正概率类之和必须在 $10^{-12}$ 容差内为一。零概率类不输出，从而减少 JSON 体积；这不改变概率总量。每一类保留
环境名、信息元件状态、功能结果、保护清除类型、失效原因和 DER-FRT 终态。

\subsection{三窗口后果}
联合类的事件能量由清除、切换和修复三个连续窗口组成：
\begin{align}
 E_c={}&S_c^{clear}\frac{t_c^{term}}{3600}
 +S_c^{switch}\tau_c^{switch}
 +S_c^{residual}\max(0,\tau^{repair}-\tau_c^{switch}).
\end{align}
主清除、后备清除和未清除分别选择不同清除切负荷；隔离失败选择扩大区切负荷；恢复失败保留残余切负荷。DER 终态造成的
发电损失等值切负荷逐设备相加，并限制恢复比例到 $[0,1]$。

\subsection{序贯 CTMC 元件过程}
序贯信息物理模拟为每个信息元件建立故障/修复事件。若年故障率为 $\lambda_q$、平均修复时间为 $r_q$，等待时间分别为
\begin{equation}
 T_q^{up}\sim\operatorname{Exp}(\lambda_q/H),\qquad
 T_q^{down}\sim\operatorname{Exp}(1/r_q).
\end{equation}
事件队列使用连续小时，保留亚小时动作，不向上取整到整点。相同时间事件用单调序号稳定排序，保证固定随机种子下执行顺序
确定。

\subsection{共因事件过程}
共因组以发生率和平均持续时间生成开始/结束事件。开始时组内全部元件强制失效，结束后恢复到各自固有状态；若个别元件
本身仍处于故障修复期，则共因结束不会错误地把它置为可用。结果返回共因事件数和
\fld{common\_cause\_modelled}，便于区分“输入为空”和“模型未消费”。

\subsection{物理故障和供电恢复事件}
物理故障可以按泊松过程生成，也可提供确定发生时刻用于解析回归。故障使绑定母线失电并安排供电恢复事件；其间信息节点
从备用电池取能。供电恢复只解除物理失电约束，不会修复固有信息故障或共因故障。

\subsection{电池连续耗尽}
在事件间隔 $[t_k,t_{k+1})$ 内，失电信息节点能量按
\begin{equation}
 E_q(t_{k+1})=\max\{0,E_q(t_k)-P_q^{draw}(t_{k+1}-t_k)\}
\end{equation}
更新。若耗尽时刻落在区间内部，模拟器插入耗尽边界而不是等到下一外部事件。这保证保护功能在电池耗尽前后可以选择不同
分支，并返回全程最小电池能量。

\subsection{包交付的序贯抽样}
静态枚举把 $A_qd_q$ 合成条件成功概率；序贯路径在每次物理故障需求时对包交付作伯努利抽样，同时保留元件当前健康、
QoS 门限和共享路径。结果标志 \fld{packet\_delivery\_sampled} 只有实际执行过需求抽样时才为真，不能由配置字段预置。

\subsection{保护类抽样}
给定当前确定信息状态，模拟器生成条件保护类，再按类概率抽取主清除、后备清除或未清除。机械成功概率在每次需求上重新
抽样；不会把断路器某次拒动永久保留为元件故障，除非调用方另建一个具有修复过程的断路器状态元件。

\subsection{停电区间并集}
多个故障可能产生重叠停电。简单累加每个事件的 $S\tau$ 会重复计算相同负荷和时段。当前序贯接口对同一代表性切负荷强度
形成区间集合，排序后合并：
\begin{equation}
 \mathcal I=\bigcup_k[t_k^{start},t_k^{end}).
\end{equation}
LOLE 取并集长度，LOLF 取不连续合并区间数，EENS 取区间长度乘切负荷。结果返回合并区间数和
\fld{outage\_interval\_union\_modelled}。若同时需要多个不同负荷点，应扩展为逐负荷点区间并集，不能用系统最大切负荷代替。

\subsection{静态与序贯结果的差异来源}
静态事件树用环境概率和代表时长直接求期望，适合参数筛选；序贯队列显式表示故障先后、电池跨事件状态、重叠停电和共因
持续过程。两者在低频、无重叠、每次事件前电池充满且持续时间固定时应趋近；电池未充满、事件重叠或共因持续时，序贯结果
不应被强求等于静态结果。

\subsection{边界和复杂度}
静态元件位枚举复杂度 $O(2^{n_c}n_fn_p)$；序贯队列主要复杂度为 $O(N_e\log N_e+N_kn_p)$，其中 $N_e$ 是事件数、
$N_k$ 是物理故障需求数。\fld{maximum\_events} 防止异常高率输入导致无界运行；达到上限明确报错，不返回不完整年度均值。

\subsection{解析回归算例}
测试安排两个确定物理故障，其中第二次发生在信息节点备用电池耗尽之后。第一次自动保护成功，第二次因信息功能丢失进入
保守后果。预期 EENS 为 $10(1+1/3600)$ MWh，LOLE 为 $1+1/3600$ 小时；该非整数小时结果证明事件时标未被整点化。
另一个测试构造重叠停电，验证 LOLE 取区间并集而非各事件持续时间之和。

\subsection{输出真实性标志}
序贯结果分别报告亚小时计时、物理供电耦合、共因、保护概率、包交付、电池轨迹和区间并集标志。标志由实际执行路径设置：
空共因数组不能得到“共因已建模”，没有电池节点不能得到“电池轨迹已建模”。该约束防止 GUI 仅因用户勾选选项就显示
并未消费的数据能力。
